This paper documents a fundamental transformation in broadband demand across 33 European countries over 2010--2024. Three findings emerge from our two-way fixed effects analysis with Driscoll--Kraay standard errors.

First, substantial regional heterogeneity characterised pre-COVID demand: EaP countries exhibited strong price sensitivity ($\varepsilon = \EaPBaseBr$, p$<$0.001)---\EaPEURatio{} times larger than EU elasticity ($\varepsilon = \EUBaseBr$, p$<$0.10)---reflecting income-driven affordability constraints. Second, both regions converged to near-zero elasticity during 2020--2024, with price changes having no detectable effect on adoption---evidence that demand-side digital harmonisation between EaP and EU markets has already been achieved. Third, and most critically, the pre-COVID trend is detectable before 2020: for EU countries, year-by-year estimates show elasticity declining from $\EUYbyYRefElast$ in 2015 toward zero; for EaP countries, the pooled placebo test confirms the same direction (p$=$\PlaceboEaPPval), as EaP's small sample (five to six countries) precludes reliable annual estimates. Together, these findings indicate the pandemic accelerated but did not initiate broadband's integration into economic and social life. Methodologically, we demonstrate that price measurement critically affects inference---income-relative prices yield robust results while PPP-adjusted prices do not.

For policy, the near-zero elasticity implies that affordability interventions---effective when demand was elastic throughout the pre-COVID period---have lost traction. Policy must now prioritize infrastructure deployment and universal service obligations. Methodologically, the finding that only income-relative prices yield robust elasticity estimates implies that policymakers tracking broadband affordability should use GNI-based metrics rather than PPP-adjusted or nominal prices. Broadband increasingly resembles utilities like water or electricity, potentially justifying stronger regulatory oversight and public investment to ensure equitable access.

Future research should examine household-level heterogeneity to identify remaining price-sensitive populations, extend analysis to mobile broadband and other regions, and develop quality-adjusted price indices. As broadband completes its transformation into essential infrastructure, understanding these dynamics becomes crucial for bridging digital divides worldwide.
